\section{Experimental Setup}
\label{sec:setup}

\paragraph{Task grid.}
We evaluate five optimizer models from three developers on all four optimization tasks in our suite:
\texttt{claude-opus-5}, \texttt{claude-sonnet-5}, \texttt{gpt-5.6-sol},
\texttt{gpt-5.6-terra}, and \texttt{kimi-k3}, chosen as the current frontier
release of each family. Each model is paired with two coding harnesses: a
\textbf{shared harness} (\oc), held fixed across all models, and the model's
\textbf{native harness}, namely \cc for the Claude models, \codex for the GPT
models, and \kc for Kimi. We refer to each model--harness pair as an
\emph{optimizer configuration}, giving \NumGridOptimizers\ configurations in our core grid. Comparing models under \oc holds the coding harness fixed; comparing a model
across \oc and its native harness measures sensitivity to scaffold choice. 
We run two additional harnesses -- \goose and \miniswe -- across all models on GAIA to compare sensitivity across optimizer harnesses. Finally, for two of the five models -- \texttt{claude-opus} and \texttt{gpt-5.x} -- we run earlier releases of the same family using each family's native harness on OfficeQA,
supporting the capability ladder we present in Section~\ref{sec:ladder}.
The complete per-task results, including observed run ranges and properties of
the generated harnesses, are reported in Appendix~\ref{app:master}
(Table~\ref{tab:master}).

\paragraph{Scoring protocol.}
Scoring follows the protocol of Section~\ref{sec:scoring}: each held-out evaluation is the
mean of three attempts per test case, matching the $K{=}3$ pooling behind each
task's pinned baseline. Each optimizer configuration is run twice; we report
ranges in Table~\ref{tab:master}. The optimizer model's inference is metered but left
uncapped, so these results estimate what is achievable when the optimizer model's
reasoning is not the scarce resource.

\paragraph{Analysis protocol.}
\label{sec:analysis}
Three choices govern the analyses in Section~\ref{sec:results}.

\begin{itemize}[itemsep=2pt,topsep=2pt,leftmargin=*]

  \item \textbf{Outcome.} We report the normalized gain $g$ of
  Eq.~\ref{eq:normalized-gain}: the fraction of the headroom above the pinned
  baseline that an optimizer captured.

  \item \textbf{Measurement resolution.} We estimate evaluation noise by
  scoring the same candidate twice on the same cases and carry the median
  discrepancy to the $K{=}3$ normalized-gain scale used for held-out scoring.
  The resulting task-specific \emph{resolution band} is a descriptive threshold:
  differences smaller than the band are treated as unresolved, not as formal
  significance-test results. Table~\ref{tab:gain-by-task} reports each task's
  band.

  \item \textbf{Composite model score.} Normalized gain is defined within a
task; normalization places tasks in common headroom units but does not remove
systematic differences among them. To obtain a cross-task score for optimizer
model $m$, let $g_{mtr}$ denote its normalized gain on task $t$ in qualifying
replicate $r$ under the shared harness, and first average replicates:
\begin{equation}
\bar{g}_{mt}
=
\frac{1}{|\mathcal{R}_{mt}|}
\sum_{r\in\mathcal{R}_{mt}} g_{mtr}.
\end{equation}
We then decompose these configuration-level gains as
\begin{align*}
\bar{g}_{mt}
=
\mu+\tau_t+\lambda_m+\varepsilon_{mt},
\qquad \\
\sum_t\tau_t=0,
\quad
\sum_m\lambda_m=0,
\label{eq:lss}
\end{align*}
where $\tau_t$ absorbs task-level differences and $\lambda_m$ is the optimizer
model effect. Because the shared-harness grid is balanced,
\begin{equation}
\widehat{\lambda}_m
=
\frac{1}{|\mathcal{T}|}\sum_{t\in\mathcal{T}}\bar{g}_{mt}
-
\frac{1}{|\mathcal{M}||\mathcal{T}|}
\sum_{m'\in\mathcal{M}}\sum_{t\in\mathcal{T}}\bar{g}_{m't}.
\end{equation}
We define
$\operatorname{LSS}_{\lambda}(m)=\widehat{\lambda}_m$. Thus
LSS-$\lambda$ is the model's task-adjusted mean performance relative to the
evaluated models' grand mean, expressed in normalized-gain units; higher is
better. The primary score uses the shared-harness runs on the tasks with
competent seeds. Native-harness runs are excluded to hold the scaffold fixed.

\end{itemize}

% Declared one unit earlier than it is discussed, on purpose. A full-width float
% cannot be placed on the page it is declared on -- LaTeX's output routine only
% considers double floats at a page boundary -- so a figure* declared partway down
% page N is deferred to the top of page N+1. Moving the declaration back is the only
% way to pull the figure forward; stfloats/dblfloatfix add bottom placement on the
% later page and do not help with this.
\begin{figure*}[t]
\centering
\includegraphics[width=\linewidth]{figures/R03_levers_vs_gain.pdf}
\caption{\textbf{Explored breadth is associated with gain.} Normalized gain
against the fraction of eight pre-specified harness levers touched during
search. Each point averages replicate runs for one optimizer configuration;
Spearman $\rho$ is computed separately by task. Lever coverage measures
exploration, not changes retained in the submitted harness, and is correlated
with overall modification volume.}
\label{fig:levers}
\end{figure*}
